%-----------------------------------------------------------------------------------------
% Autor dieser Vorlage:
% Stefan Macke (http://fachinformatiker-anwendungsentwicklung.net)
% Permalink zur Vorlage: http://fiae.link/LaTeXVorlageFIAE
%
% Sämtliche verwendeten Abbildungen, Tabellen und Listings stammen von Dirk Grashorn.
%
% Lizenz: Creative Commons 4.0 Namensnennung - Weitergabe unter gleichen Bedingungen
% -----------------------------------------------------------------------------------------

\documentclass[
	ngerman,
	toc=listof, % Abbildungsverzeichnis sowie Tabellenverzeichnis in das Inhaltsverzeichnis aufnehmen
	toc=bibliography, % Literaturverzeichnis in das Inhaltsverzeichnis aufnehmen
	footnotes=multiple, % Trennen von direkt aufeinander folgenden Fußnoten
	parskip=half, % vertikalen Abstand zwischen Absätzen verwenden anstatt horizontale Einrückung von Folgeabsätzen
	numbers=noendperiod % Den letzten Punkt nach einer Nummerierung entfernen (nach DIN 5008)
]{scrartcl}
\pdfminorversion=5 % erlaubt das Einfügen von pdf-Dateien bis Version 1.7, ohne eine Fehlermeldung zu werfen (keine Garantie für fehlerfreies Einbetten!)
\usepackage[utf8]{inputenc} % muss als erstes eingebunden werden, da Meta/Packages ggfs. Sonderzeichen enthalten

% !TEX root = Projektdokumentation.tex

% Hinweis: der Titel muss zum Inhalt des Projekts passen und den zentralen Inhalt des Projekts deutlich herausstellen
\newcommand{\titel}{Entwicklung von NatInfo}
\newcommand{\untertitel}{Webbasiertes Tool zur Unterstützung der Entwickler}
\newcommand{\kompletterTitel}{\titel{} -- \untertitel}

\newcommand{\autorName}{Stefan Macke}
\newcommand{\autorAnschrift}{Meine Straße 1}
\newcommand{\autorOrt}{49377 Vechta}

\newcommand{\betriebLogo}{LogoBetrieb.pdf}
\newcommand{\betriebName}{\textsc{Alte Oldenburger} Krankenversicherung AG}
\newcommand{\betriebAnschrift}{Alte-Oldenburger-Platz 1}
\newcommand{\betriebOrt}{49377 Vechta}

\newcommand{\ausbildungsberuf}{Fachinformatiker für Anwendungsentwicklung}
\newcommand{\betreff}{Dokumentation zur betrieblichen Projektarbeit}
\newcommand{\pruefungstermin}{Sommer 2015}
\newcommand{\abgabeOrt}{Vechta}
\newcommand{\abgabeTermin}{23.04.2015}
 % Metadaten zu diesem Dokument (Autor usw.)
\input{Allgemein/Packages} % verwendete Packages
\input{Allgemein/Seitenstil} % Definitionen zum Aussehen der Seiten
% Abkürzungen
\newcommand{\Versis}{\textsc{Versis}\xspace}
\newcommand{\NI}{NatInfo\xspace}
\newcommand{\AO}{\textsc{Alte Oldenburger} Krankenversicherung\xspace}
 % eigene allgemeine Befehle, die z.B. die Arbeit mit LaTeX erleichtern
% Abkürzungen
\newcommand{\Versis}{\textsc{Versis}\xspace}
\newcommand{\NI}{NatInfo\xspace}
\newcommand{\AO}{\textsc{Alte Oldenburger} Krankenversicherung\xspace}
 % eigene projektspezifische Befehle, z.B. Abkürzungen usw.

\begin{document}

\phantomsection
\thispagestyle{empty}
\pdfbookmark[1]{Eidesstattliche Erklärung}{ihkdeckblatt}
\includegraphicsKeepAspectRatio{DeckblattIHK}{1}
\cleardoublepage

\phantomsection
\thispagestyle{plain}
\pdfbookmark[1]{Deckblatt}{deckblatt}
% !TEX root = Projektdokumentation.tex
\begin{titlepage}

    %TODO: IHK Logo ersetzen
\begin{center}
\includegraphics[scale=0.25]{LogoIHK.pdf}\\[1ex]
\Large{Abschlussprüfung \pruefungstermin}\\[3ex]

\Large{\ausbildungsberuf}\\
\LARGE{\betreff}\\[4ex]

\huge{\textbf{\titel}}\\[1.5ex]
\Large{\textbf{\untertitel}}\\[4ex]

\normalsize
Abgabetermin: \abgabeOrt, den \abgabeTermin\\[3em]
\textbf{Prüfungsbewerber:}\\
\autorName\\
\autorAnschrift\\
\autorOrt\\[5ex]

\includegraphics[scale=0.4]{\betriebLogo}\\[2ex]
\textbf{Ausbildungsbetrieb:}\\
\betriebName\\
\betriebAnschrift\\
\betriebOrt\\[5em]
\end{center}

\end{titlepage}
\cleardoublepage

% Preface --------------------------------------------------------------------
\phantomsection
\pagenumbering{Roman}
\pdfbookmark[1]{Inhaltsverzeichnis}{inhalt}
\tableofcontents

\cleardoublepage

\phantomsection
\listoffigures
\cleardoublepage

\phantomsection
\listoftables
\cleardoublepage

\phantomsection
\lstlistoflistings
\cleardoublepage

\newcommand{\abkvz}{Abkürzungsverzeichnis}
\renewcommand{\nomname}{\abkvz}
\section*{\abkvz}
\markboth{\abkvz}{\abkvz}
\addcontentsline{toc}{section}{\abkvz}
% !TEX root = Projektdokumentation.tex

% Es werden nur die Abkürzungen aufgelistet, die mit \ac definiert und auch benutzt wurden. 
%
% \acro{VERSIS}{Versicherungsinformationssystem\acroextra{ (Bestandsführungssystem)}}
% Ergibt in der Liste: VERSIS Versicherungsinformationssystem (Bestandsführungssystem)
% Im Text aber: \ac{VERSIS} -> Versicherungsinformationssystem (VERSIS)

% Hinweis: allgemein bekannte Abkürzungen wie z.B. bzw. u.a. müssen nicht ins Abkürzungsverzeichnis aufgenommen werden
% Hinweis: allgemein bekannte IT-Begriffe wie Datenbank oder Programmiersprache müssen nicht erläutert werden,
%          aber ggfs. Fachbegriffe aus der Domäne des Prüflings (z.B. Versicherung)

% Die Option (in den eckigen Klammern) enthält das längste Label oder
% einen Platzhalter der die Breite der linken Spalte bestimmt.
\begin{acronym}[WWWWW]
	\acro{AO}{\textsc{Alte Oldenburger} Krankenversicherung AG}
	\acro{API}{Application Programming Interface}
	\acro{CSS}{Cascading Style Sheets}
	\acro{CSV}{Comma Separated Value}
	\acro{EPK}{Ereignisgesteuerte Prozesskette}
	\acro{ERM}{En\-ti\-ty-Re\-la\-tion\-ship-Mo\-dell}
	\acro{HTML}{Hypertext Markup Language}\acused{HTML}
	\acro{IDE}{Integrated Development Environment}
	\acro{MVC}[MVC]{Model View Controller}
	\acro{NatInfo}[\textsc{NatInfo}]{Natural Information System}
	\acro{Natural}[\textsc{Natural}]{Programmiersprache der Software AG}\acused{Natural}
	\acro{ORM}{Object-Relational Mapping}
	\acro{PHP}{Hypertext Preprocessor}
	\acro{SDK}{Software Development Kit}
	\acro{SQL}{Structured Query Language}
	\acro{SVN}{Subversion}
	\acro{UML}{Unified Modeling Language}
	\acro{XML}{Extensible Markup Language}
\end{acronym}

\clearpage

% Inhalt ---------------------------------------------------------------------
\pagenumbering{arabic}
% !TEX root = Projektdokumentation.tex
%! Author = apena
%! Date = 30.06.24

\section{Ausgangssituation}
\label{sec:ausgangssituation}

\subsection{Projektziele}
\label{subsec:projektziele}

\subsection{Teilaufgaben}
\label{subsec:teilaufgaben}

\subsection{Kundenwünsche}
\label{subsec:kundenwuensche}

\subsection{Projektumfeld}
\label{subsec:projektumfeld}

\subsection{Prozessschnittstellen}
\label{subsec:prozessschnittstellen}
%! Author = apena
%! Date = 30.06.24

\section{Ressourcenplanung}
\label{sec:ressourcenplanung}

\subsection{Personalplanung}
\label{subsec:personalplanung}

\subsection{Sachmittelplanung}
\label{subsec:sachmittelplanung}

\subsection{Terminplanung}
\label{subsec:terminplanung}

\subsection{Kostenplanung}
\label{subsec:kostenplanung}


%! Author = apena
%! Date = 30.06.24

\section{Ablaufplanung}
\label{sec:ablaufplanung}


%! Author = apena
%! Date = 30.06.24

\section{Durchführung und Auftragsbearbeitung}
\label{sec:durchfuehrung}

\subsection{Prozessschritte}
\label{subsec:prozessschritte}

\subsection{Vorgehensweise}
\label{subsec:vorgehensweise}

\subsection{Qualitätssicherung}
\label{subsec:qualitaetssicherung}

\subsection{Abweichungen}
\label{subsec:abweichungen}

\subsection{Anpassungen}
\label{subsec:anpassungen}

\subsection{Entscheidungen}
\label{subsec:entscheidungen}
%! Author = apena
%! Date = 30.06.24

\section{Projektergebnisse}
\label{sec:projektergebnisse}

\subsection{Soll-Ist-Vergleich}
\label{subsec:soll-ist-vergleich}

\subsection{Bewertung}
\label{subsec:bewertung}

% Literatur ------------------------------------------------------------------
\clearpage
\renewcommand{\refname}{Literaturverzeichnis}
\bibliography{Bibliographie}
\bibliographystyle{Allgemein/natdin} % DIN-Stil des Literaturverzeichnisses
\input{Erklaerung}

% Anhang ---------------------------------------------------------------------
\clearpage
\appendix
\pagenumbering{roman}
% !TEX root = Projektdokumentation.tex
\section{Anhang}
\subsection{Detaillierte Zeitplanung}
\label{app:Zeitplanung}

\tabelleAnhang{ZeitplanungKomplett}

\input{Anhang/AnhangLastenheft.tex}
\clearpage

\subsection{Use Case-Diagramm}
\label{app:UseCase}
Use Case-Diagramme und weitere \acs{UML}-Diagramme kann man auch direkt mit \LaTeX{} zeichnen, siehe \zB \url{http://metauml.sourceforge.net/old/usecase-diagram.html}.
\begin{figure}[htb]
\centering
\includegraphicsKeepAspectRatio{UseCase.pdf}{0.7}
\caption{Use Case-Diagramm}
\end{figure}

\input{Anhang/AnhangPflichtenheft.tex}

\subsection{Datenbankmodell}
\label{app:Datenbankmodell}
ER-Modelle kann man auch direkt mit \LaTeX{} zeichnen, siehe \zB \url{http://www.texample.net/tikz/examples/entity-relationship-diagram/}.
\begin{figure}[htb]
\centering
\includegraphicsKeepAspectRatio{database.pdf}{1}
\caption{Datenbankmodell}
\end{figure}
\clearpage

\input{Anhang/AnhangEntwuerfe.tex}
\clearpage
\input{Anhang/AnhangScreenshots.tex}
\input{Anhang/AnhangDoc.tex}
\clearpage
\subsection{Testfall und sein Aufruf auf der Konsole}
\label{app:Test}
\lstinputlisting[language=php, caption={Testfall in PHP}]{Listings/Out/tests.php}
\clearpage
\begin{figure}[htb]
\centering
\includegraphicsKeepAspectRatio{testcase.jpg}{1}
\caption{Aufruf des Testfalls auf der Konsole}
\end{figure}


\subsection{Klasse: ComparedNaturalModuleInformation}
\label{app:CNMI}
Kommentare und simple Getter/Setter werden nicht angezeigt.
\lstinputlisting[language=php, caption={Klasse: ComparedNaturalModuleInformation}]{Listings/Out/cnmi.php}
\clearpage

\subsection{Klassendiagramm}
\label{app:Klassendiagramm}
Klassendiagramme und weitere \acs{UML}-Diagramme kann man auch direkt mit \LaTeX{} zeichnen, siehe \zB \url{http://metauml.sourceforge.net/old/class-diagram.html}.
\begin{figure}[htb]
\centering
\includegraphicsKeepAspectRatio{Klassendiagramm.pdf}{1}
\caption{Klassendiagramm}
\end{figure}
\clearpage

\input{Anhang/AnhangBenutzerDoku.tex}


\end{document}
